% Generated draft: requires mathematical and source review.
\documentclass[11pt]{article}
\usepackage[margin=25mm]{geometry}
\usepackage{fontspec}
\setmainfont{Latin Modern Roman}
\usepackage{amsmath,amssymb,mathtools,unicode-math}
\setmathfont{Latin Modern Math}
\usepackage{enumitem,needspace,xurl,hyperref}
\setlength{\parindent}{0pt}
\setlength{\parskip}{5pt}
\setlength{\emergencystretch}{3em}
\newcommand{\N}{\mathbb{N}}
\newcommand{\Z}{\mathbb{Z}}
\newcommand{\Q}{\mathbb{Q}}
\newcommand{\R}{\mathbb{R}}
\newcommand{\C}{\mathbb{C}}
\newcommand{\F}{\mathbb{F}}
\newcommand{\A}{\mathbb{A}}
\newcommand{\PP}{\mathbb{P}}
\newcommand{\E}{\mathbb{E}}
\newcommand{\eps}{\varepsilon}
\newcommand{\dd}{\,\mathrm{d}}
\begin{document}
\section*{UnsolvedMath ID: 30004871}
\textbf{Source problem code:} catalogue entry
\subsection{Definitions and mathematical statement}
For purposes of interpreting arithmeticity and thinness, suppose an integral realization of a finitely generated group \(\Gamma\leq \mathrm{GL}_n(\Z)\) has been specified. Let \(G=\overline{\Gamma}^{\,\mathrm{Zar},\Q}\) be its algebraic Zariski closure over \(\Q\), and put
\[
G(\Z)=G(\Q)\cap\mathrm{GL}_n(\Z).
\]
In this realization, \(\Gamma\) is called arithmetic if \([G(\Z):\Gamma]<\infty\), and thin if it is Zariski dense in \(G\) and \([G(\Z):\Gamma]=\infty\). The supplied excerpt does not specify the integral realization or the admissible hypergeometric parameters \(\alpha,\beta\); these definitions therefore explain the terminology conditionally and do not supply the missing parameter restrictions.

The source calls a hypergeometric monodromy group \(\Gamma_{\alpha,\beta}\) a \emph{hyperbolic hypergeometric monodromy group} (HHMG) when its associated algebraic group \(G_{\alpha,\beta}\) is of special orthogonal type \(\mathrm{SO}_n\) and
\[
G_{\alpha,\beta}(\R)\cong\mathrm{SO}(n-1,1).
\]
This is the source's stated real-group condition; the relevant rational quadratic form and the construction of \(G_{\alpha,\beta}\) are not supplied.

\textbf{Classification question.} Among the HHMGs in the source's intended hypergeometric family, determine which \(\Gamma_{\alpha,\beta}\) are arithmetic and which are thin.

\textbf{Conjecture attributed to Fuchs--Meiri--Sarnak (Conjecture 9).} All but finitely many HHMGs are thin. Formally, if \(\mathcal H\) denotes the intended collection of HHMGs with the source's intended convention for identifying groups, the assertion is
\[
\#\{\Gamma\in\mathcal H:\Gamma\text{ is not thin}\}<\infty.
\]
Neither the complete definition of \(\mathcal H\) nor the convention for counting distinct groups is included in the supplied record. Consequently, this quantified formulation remains conditional on those missing specifications.

The surrounding Question 8 asks the arithmeticity-versus-thinness question for the broader family of hypergeometric monodromy groups. The excerpt reports arithmeticity for \(n\in\{2,3\}\), further work in the symplectic case \(G_{\alpha,\beta}=\mathrm{Sp}_n\), and seven families together with sporadic examples in the hyperbolic case. It attributes to Fuchs--Meiri--Sarnak a certificate for thinness and proofs of thinness for some of those families and many sporadic examples. These are supplied contextual claims, subject to the omitted hypotheses of the underlying family, rather than independently verified assertions here.
\subsection{Short English statement}
Classify the arithmetic and thin hyperbolic hypergeometric monodromy groups, and formulate the Fuchs--Meiri--Sarnak conjecture that only finitely many are not thin. The supplied record omits the admissible parameters, integral realizations, and convention for counting distinct groups, so a fully specified quantification requires further source review. The supplied open status is dated curation.
\subsection{Sources}
\begin{itemize}
\item Supplied local UnsolvedMath registry record, permanent ID \(30004871\), exact source problem code \texttt{catalogue entry}, in the Oberwolfach Reports Open Problems collection. The record labels the problem \texttt{open} and the statement \texttt{corrected\_verified}; those are supplied registry labels, not verification performed for this draft.
\item \emph{Arbeitsgemeinschaft: Thin Groups and Super-approximation} (2022), \emph{Oberwolfach Reports}, pp.\ 2721--2790, DOI \texttt{10.4171/owr/2021/50}. Supplied authoritative report reference: \url{https://doi.org/10.4171/owr/2021/50}. The extracted passage identifies Question 8 and Conjecture 9 and includes the printed page marker 2739. Conjecture 9 is attributed to Fuchs--Meiri--Sarnak.
\item Supplied literature reference labelled ``Fuchs--Meiri--Sarnak survey'': \url{https://www.math.ucdavis.edu/~efuchs/ThinUbiquity.pdf}. No further bibliographic details are supplied.
\item Literature assessment dated exactly \texttt{2026-08-22}: the supplied curation describes the assertion as remaining a conjecture and reports families and sporadic examples rather than a general proof. This is supplied dated curation, not an independently verified literature assessment. The record also reports that the official Oberwolfach report was checked directly during its earlier statement verification; that check has not been repeated here.
\end{itemize}
\end{document}
